\section{The actual graded Gram bound for Type I}
\label{sec:graded-gram}

\status{This supplies the previously retained graded initial-Gram premise
of TI6, together with the full-word \texttt{HypP}/\texttt{HypW} inputs.
The argument controls the actual coefficient-image lattice.
It uses the proved IW averaging theorem and retains no additional local
integral representation-theory premise.}

\paragraph{GG1: the lattice and its normalization.}
Let a signed word have $K+1$ letters of degrees $d_i$, total degree
$D=\sum_i d_i$, and final vector letter $W$. Work over $\mathbb Q_p$,
with integral coefficient lattice $\mathbb Z_p[X]$.
Assume $r\ge2$, put $\rho=r-1$, and $z=W_\rho$.
For finitely many prefix bounds $(J_\ell,M_\ell)$, let $Y$ be the
multihomogeneous space whose types at prefix $J_\ell$ have first
highest weight at most $M_\ell$; $M_\ell=\infty$ is allowed.
Let $M_w$ be its integral lattice of full weight $w$, highest for
$GL_{\{0,\ldots,\rho-1\}}$ acting on all letters. The literal
\texttt{grLat} is
\[
 \begin{gathered}
 {\cal G}_e=\operatorname{in}_e(M_w\cap z^e\mathbb Q_p[X]),\qquad
 \operatorname{in}_e\text{ retains }z^e,\\
 B_d(X^\alpha,X^\beta)=
 \delta_{\alpha\beta}\frac{\prod_{i,a}\alpha_{ia}!}{\prod_i d_i!}.
 \end{gathered}
                                                               \tag{GG.1}
\]
The diagonal entries are reciprocals of integer multinomials.

\paragraph{GG2: all required integral projectors.}
Put
\[
 e_1=\binom r2\operatorname{Nat.log}_p(D+r-1),\qquad
 D_4=\binom r2\operatorname{Nat.log}_p(4D+r-1).                  \tag{GG.2}
\]
Any two prefix weights differ at one coordinate by at most $D$:
each lies between minus the covector degree and the vector degree.
Elementary substitutions show that the coefficient lattice is
weight-graded and stable under root divided powers. IW4 therefore
makes every prefix type projector integral after multiplication by
$p^{e_1}$. They are Fock self-adjoint and commute for nested prefixes,
because the smaller projector commutes with the larger diagonal action.

On a fixed full weight $w$, the smaller-group projector onto type
$w|_{\{0,\ldots,\rho-1\}}$ is precisely the highest-vector projector:
the top-weight space consists of highest vectors, and every other
type contributes lowering images. This proves \texttt{HypE} with the
same cost.

For an isotypic lattice $M$, let $\Delta$ be the divided-power span
of its integral highest vectors. Its integral top-weight projector
$C$ satisfies, by IW2 on $\operatorname{Hom}(M,\Delta)$,
$I(C)=(\dim U)^{-1}\operatorname{Id}$, where $U$ is one irreducible
factor. Multiplying by the integer $\dim U$ gives
$p^{e_1}M\subset\Delta$. The multiplicity space is untouched.
This proves \texttt{HypW} without integral semisimplicity.

\paragraph{GG3--GG4: the actual restriction-image lattice.}
Allow the final degree to be $o$. Let $X_o$ consist of $z$-free
polynomials satisfying the earlier prefix bounds. Each full-prefix
bound becomes both the earlier-word $GL_r$ bound and the full-word
$GL_{r-1}$ bound. These are exactly \texttt{XSp}.
If all full bounds are at least $D$, restriction lifts trivially.
Otherwise let $c<D$ be their finite minimum and set
\[
 {\cal W}=\bigoplus_oX_o,\qquad
 {\cal L}_{\rm im}=\sum_o
       \{h|_{z=0}:h\in Y_o\cap\mathbb Z_p[X]\}.                 \tag{GG.3}
\]
Restriction lands in $X_o$: extracting $W$-coefficients tensors
with $\operatorname{Sym}^o(V)^*$, whose coordinate weights are
nonpositive, and hence preserves the earlier first-weight bound.

Write $A_{ab}$ for the earlier action,
$L_{ab}=A_{ab}+W_a\partial_b$, and $E_W=\sum W_i\partial_i$.
On the remaining coordinates define
\[
 U_a=W_a(c-E_W)-\sum_bW_bA_{ab},\qquad
 E_{**}=c-E_W,\quad E_{*b}=\partial_b,\quad
 E_{a*}=U_a,\quad E_{ab}=L_{ab}.                               \tag{GG.4}
\]
Direct commutators give $\mathfrak{gl}_r$.
The family $U_a=cW_a-\sum_bW_bL_{ab}$ is covariant.
On coordinate-$0$ weight $c$, $U_0$ vanishes; the other $U_a$
do not raise that weight. Thus it preserves the first-weight cutoff.
Derivatives form a covector family and preserve it as well.
This proves stability directly, without the source's intermediate
rational-surjectivity theorem.

Every nonzero $X_o$ has $o\le c+D_{\rm early}$: a generic coordinate
change makes its $W_0^o$ coefficient nonzero, whose earlier signed
weight is at least $-D_{\rm early}$. The auxiliary weights therefore
lie in $[-D_{\rm early},c+2D_{\rm early}]$, of width at most $4D$.

Divided powers are integral. For $U_a$, put $q=1+tW_a$.
Its exponential is $q^c$ times substitution
\[
 W_j\mapsto W_j/q,\quad
 x_b\mapsto x_b-tW_bx_a/q,\quad
 y_a\mapsto y_a+t\sum_bW_by_b,                                 \tag{GG.5}
\]
on earlier vector and covector letters; other covector coordinates
are fixed. Integer powers of $q$ have integral series coefficients.
The same formula before restriction proves that $ {\cal L}_{\rm im}$
is divided-power stable.

An auxiliary highest vector is killed by every $\partial_b$, so has
$W$-degree zero and lifts to itself. IW4 first isolates each type of
an integral vector in $ {\cal W}$, at cost $D_4$; GG2's highest-vector
generation costs another $D_4$. Their generators lie in
$ {\cal L}_{\rm im}$. Taking the original degree component proves
\[
 p^{2D_4}(X_{d_K}\cap\mathbb Z_p[X])
       \subset\{h|_{z=0}:h\in Y_d\cap\mathbb Z_p[X]\}.           \tag{GG.6}
\]

\paragraph{GG5--GG6: order, saturation and projection.}
Dividing by $z^e$ reduces the final degree and full weight by $e$,
and subtracts $e$ from precisely the full-prefix bounds.
This is an equivalence: after a generic group change, multiplication
by $z^e$ adds $e$ to the largest first-coordinate Laurent exponent.
Project the lift in (GG.6) to its full weight and smaller-group
highest vectors, costing $e_1$. If $V_e=\mathbb Q_p{\cal G}_e$,
\[
 p^b(V_e\cap\mathbb Z_p[X])\subset{\cal G}_e,\qquad
 b=e_1+2D_4.                                                   \tag{GG.7}
\]
This establishes saturation control for the actual image.

For $K\ge1$, first project integrally to $z$-degree zero. The rational
projection onto the reduced highest restriction space then uses
$K+2$ nested type projectors. Those at prefix
lengths $0,1$ are identity or zero. The remaining $K$ cost $Ke_1$;
highest projection costs $e_1$. Strip $z^e$, apply this product and
multiply back. It gives a Fock-orthogonal projection $P$ onto $V_e$
with integral cost $A=(K+1)e_1$. The identity
$\operatorname{Fock}(z^ef,z^eg)=e!\operatorname{Fock}(f,g)$
preserves orthogonality.

\paragraph{GG7--GG8: determinant and the precise segment constant.}
For primitive $u\in{\cal G}_e$, its smallest coefficient valuation
is at most $b$ by (GG.7). Choose its corresponding monomial $m$.
Then $h=p^{A+b}P(m)\in{\cal G}_e$, and (GG.1) gives
$v_pB_d(u,h)\le A+2b$.
Smith form bounds every elementary-divisor valuation by this number.
Thus for any basis of rank $k$,
\[
 v_p\det B_d\le\bigl((K+3)e_1+4D_4\bigr)k
 \le (K+7)\frac{r^2}{2}\ell\,k,                                \tag{GG.8}
\]
when both logarithms in (GG.2) are at most $\ell$.
For $K=0$ the possible graded lattice is a monomial line with Gram
entry $1/\binom{d_0}{e}$, so its valuation is at most zero.

In the actual segment let $\ell=1+\log_p(8a+r)$, $a\ge1$, $p\ge3$.
The literal size hypothesis gives full total at most $8a+r$ and
$2D_i+1\le8a+r$ for each tie-prefix total.
Both required logarithms are therefore at most $\ell$.
Since $K=\operatorname{pos}_i<2s+1$, (GG.8) proves exactly
\[
 \kappa_4=(s+4)r^2\ell.                                       \tag{GG.9}
\]
GG2 also supplies full-word \texttt{HypP}/\texttt{HypW} with the
permitted larger exponent $\lfloor r^2\ell\rfloor$ required by
the partner construction.
